\documentclass[11pt]{article}

\usepackage{amsmath,amssymb,amsthm,mathtools}

\newtheorem{theorem}{Theorem}
\newtheorem{proposition}{Proposition}
\newtheorem{lemma}{Lemma}
\newtheorem{corollary}{Corollary}
\newtheorem{remark}{Remark}

\newcommand{\C}{\mathbb C}
\newcommand{\SL}{\mathrm{SL}}
\newcommand{\GL}{\mathrm{GL}}
\newcommand{\Sym}{\operatorname{Sym}}
\newcommand{\Schur}{\mathbb S}
\newcommand{\ch}{\operatorname{ch}}

\begin{document}

\title{Equivalence Between the Plethystic Coefficient and the \(SL_9\)-Invariant Dimension}
\author{}
\date{}
\maketitle

\begin{center}
\textbf{Status: Proved-Alg for the equivalence theorem}
\end{center}

Let \(V=\C^9\), with its standard action of \(GL(V)\simeq GL_9(\C)\), and set
\[
    M_d=\Sym^d(\Sym^3 V).
\]
This note proves that the \(SL_9\)-invariant dimension of \(M_d\) is exactly the
rectangular Schur coefficient
\[
    \big\langle h_d[h_3],s_{(k^9)}\big\rangle
\]
when \(3d=9k\), and is zero otherwise. The proof is independent of any
highest-weight rank computation.

All representations are finite-dimensional complex algebraic representations.
All Schur functors are polynomial \(GL_9\)-representations.

\section*{1. The rectangular summand and the absence of a hidden factor}

We first prove the point on which the \(GL_9\)-to-\(SL_9\) passage rests.

\begin{lemma}[Rectangular Schur functors]
Let \(V=\C^9\). For every \(k\geq 0\),
\[
    \Schur_{(k^9)}V \simeq (\det V)^k
\]
as \(GL_9\)-representations. In particular, \(\Schur_{(k^9)}V\) is
one-dimensional.
\end{lemma}

\begin{proof}
For a diagonal matrix
\[
    t=\operatorname{diag}(t_1,\ldots,t_9)\in GL_9,
\]
the character of the one-dimensional representation \((\det V)^k\) is
\[
    (\det t)^k=(t_1t_2\cdots t_9)^k=t_1^k t_2^k\cdots t_9^k.
\]
Thus \((\det V)^k\) has \(GL_9\)-weight \((k,k,\ldots,k)=(k^9)\).
Since it is one-dimensional, this weight is also its highest weight.

By the highest-weight classification of irreducible polynomial \(GL_9\)-modules,
there is a unique irreducible polynomial \(GL_9\)-module with dominant highest
weight \((k^9)\). That module is, by definition, \(\Schur_{(k^9)}V\). Hence
\[
    \Schur_{(k^9)}V\simeq(\det V)^k.
\]
The right-hand side is one-dimensional, so the same is true of
\(\Schur_{(k^9)}V\).
\end{proof}

\begin{lemma}[Fixed vectors in an irreducible polynomial \(GL_9\)-module]
Let \(\lambda=(\lambda_1,\ldots,\lambda_9)\) be a partition, with trailing
zeroes allowed. Then
\[
    \dim(\Schur_\lambda V)^{SL_9}
    =
    \begin{cases}
    1, & \lambda=(k^9)\text{ for some }k\geq 0,\\
    0, & \text{otherwise.}
    \end{cases}
\]
\end{lemma}

\begin{proof}
The restriction of a polynomial irreducible \(GL_9\)-module
\(\Schur_\lambda V\) to \(SL_9\) is the irreducible \(SL_9\)-module with highest
weight represented in simple-root coordinates by
\[
    (\lambda_1-\lambda_2,\lambda_2-\lambda_3,\ldots,\lambda_8-\lambda_9).
\]
Equivalently, one may subtract the central determinant twist:
\[
    \Schur_\lambda V
    \simeq
    (\det V)^{\lambda_9}\otimes
    \Schur_{\lambda-(\lambda_9^9)}V,
\]
and the determinant factor is trivial on \(SL_9\). The remaining highest weight
for \(SL_9\) is exactly the list of adjacent differences above.

The trivial \(SL_9\)-representation has highest weight \(0\). Hence
\(\Schur_\lambda V\) can have an \(SL_9\)-fixed vector if and only if the
restricted highest weight is zero:
\[
    \lambda_1-\lambda_2=\lambda_2-\lambda_3=\cdots=\lambda_8-\lambda_9=0.
\]
This is equivalent to
\[
    \lambda_1=\lambda_2=\cdots=\lambda_9,
\]
so \(\lambda=(k^9)\).

If \(\lambda=(k^9)\), the previous lemma gives
\[
    \Schur_{(k^9)}V\simeq(\det V)^k.
\]
For every \(g\in SL_9\), \(\det(g)=1\), so \((\det g)^k=1\). Thus this
one-dimensional representation is trivial on \(SL_9\), and its invariant
subspace has dimension \(1\).

If \(\lambda\) is not rectangular, its restricted \(SL_9\)-highest weight is
nonzero. The restricted module is then a nontrivial irreducible \(SL_9\)-module,
so it contains no copy of the trivial representation and has no nonzero
\(SL_9\)-fixed vector. Therefore
\[
    (\Schur_\lambda V)^{SL_9}=0
\]
in the non-rectangular case.
\end{proof}

\begin{remark}[No hidden multiplicity factor]
If a polynomial \(GL_9\)-module decomposes as
\[
    M\simeq\bigoplus_\lambda m_\lambda\,\Schur_\lambda V,
\]
then exactness of taking invariants over a reductive group and the preceding
lemma give
\[
    \dim M^{SL_9}
    =
    \sum_\lambda m_\lambda\,\dim(\Schur_\lambda V)^{SL_9}.
\]
Thus each copy of the rectangular summand
\(\Schur_{(k^9)}V\simeq(\det V)^k\) contributes exactly one invariant line, and
every non-rectangular summand contributes zero. Therefore
\[
    \dim M^{SL_9}=m_{(k^9)}
\]
whenever the decomposition contains a rectangular summand of total degree
\(9k\). There is no additional factor.
\end{remark}

\section*{2. Character and plethysm proof}

\begin{proposition}[Plethystic character]
The \(GL_9\)-character of
\[
    M_d=\Sym^d(\Sym^3 V)
\]
is
\[
    \ch(M_d)=h_d[h_3].
\]
If
\[
    h_d[h_3]=\sum_{\lambda\vdash 3d} a_\lambda s_\lambda,
\]
then
\[
    a_\lambda=\big\langle h_d[h_3],s_\lambda\big\rangle
\]
is the multiplicity of \(\Schur_\lambda V\) in \(M_d\), for every partition
\(\lambda\) of \(3d\) with at most \(9\) parts.
\end{proposition}

\begin{proof}
The character of \(\Sym^3V\) is \(h_3\). Applying the \(d\)-th symmetric power
to a representation corresponds, at the level of symmetric functions, to
plethysm by \(h_d\). Therefore
\[
    \ch\bigl(\Sym^d(\Sym^3V)\bigr)=h_d[h_3].
\]
The Schur functions \(s_\lambda\) are the irreducible polynomial characters of
\(GL_9\), with the convention that \(s_\lambda\) vanishes on \(9\) variables
when \(\ell(\lambda)>9\). Hence the Schur expansion
\[
    h_d[h_3]=\sum_{\lambda\vdash 3d} a_\lambda s_\lambda
\]
is exactly the isotypic decomposition
\[
    M_d\simeq
    \bigoplus_{\lambda\vdash 3d,\ \ell(\lambda)\leq 9}
    a_\lambda\,\Schur_\lambda V.
\]
Since the Schur functions are orthonormal for the Hall inner product,
\[
    a_\lambda=\big\langle h_d[h_3],s_\lambda\big\rangle.
\]
\end{proof}

\begin{theorem}[Invariant dimension equals the rectangular plethystic coefficient]
Let \(V=\C^9\) and \(M_d=\Sym^d(\Sym^3V)\). If \(3d\) is not divisible by \(9\),
then
\[
    \dim M_d^{SL_9}=0.
\]
If \(3d=9k\), then
\[
    \boxed{
    \dim M_d^{SL_9}
    =
    \big\langle h_d[h_3],s_{(k^9)}\big\rangle.
    }
\]
\end{theorem}

\begin{proof}
By the plethystic character proposition,
\[
    M_d\simeq
    \bigoplus_{\lambda\vdash 3d,\ \ell(\lambda)\leq 9}
    a_\lambda\,\Schur_\lambda V,
    \qquad
    a_\lambda=\big\langle h_d[h_3],s_\lambda\big\rangle.
\]
Taking \(SL_9\)-invariants gives
\[
    M_d^{SL_9}
    \simeq
    \bigoplus_{\lambda\vdash 3d,\ \ell(\lambda)\leq 9}
    a_\lambda\,(\Schur_\lambda V)^{SL_9}.
\]
By the fixed-vector lemma, the only Schur summands contributing
\(SL_9\)-fixed vectors are the rectangular summands
\[
    \lambda=(k,k,\ldots,k)=(k^9).
\]
Such a partition has size \(9k\), while every Schur summand in \(M_d\) has size
\(3d\). Therefore a rectangular contributor can occur only if
\[
    3d=9k.
\]
If no such integer \(k\) exists, no irreducible summand contributes any
\(SL_9\)-fixed vector, so \(\dim M_d^{SL_9}=0\).

Assume now that \(3d=9k\). The only possible contributing summand is
\(\Schur_{(k^9)}V\). By the rectangular lemma,
\[
    \Schur_{(k^9)}V\simeq(\det V)^k,
\]
which is one-dimensional and trivial on \(SL_9\). Therefore each copy of
\(\Schur_{(k^9)}V\) contributes exactly one invariant vector, and every
non-rectangular summand contributes none. Hence
\[
    \dim M_d^{SL_9}
    =
    a_{(k^9)}
    =
    \big\langle h_d[h_3],s_{(k^9)}\big\rangle.
\]
This proves the claimed equality with no extra multiplicative factor.
\end{proof}

\section*{3. Highest-weight vector formulation}

The preceding theorem identifies the invariant dimension with the plethystic
coefficient. We now record the equivalent highest-weight kernel formulation,
which computes the same multiplicity when such a rank computation is feasible.

Let \(T\subset GL_9\) be the diagonal torus and choose the standard Borel. Let
\(\varepsilon_i\) denote the \(i\)-th coordinate weight. For a \(GL_9\)-module
\(M\), write \(M_\lambda\) for the \(T\)-weight space of weight \(\lambda\), and
let
\[
    E_{i,i+1},\qquad i=1,\ldots,8,
\]
be the simple positive-root operators.

\begin{proposition}[Highest-weight kernel]
Let \(M\) be a finite-dimensional polynomial \(GL_9\)-module and let
\(\lambda\) be a dominant polynomial weight. The multiplicity of
\(\Schur_\lambda V\) in \(M\) is
\[
    \operatorname{mult}_{\Schur_\lambda V}(M)
    =
    \dim
    \bigcap_{i=1}^{8}
    \ker\left(
        E_{i,i+1}:M_\lambda\to M_{\lambda+\varepsilon_i-\varepsilon_{i+1}}
    \right).
\]
\end{proposition}

\begin{proof}
The multiplicity of \(\Schur_\lambda V\) in \(M\) equals the dimension of the
space of highest-weight vectors of weight \(\lambda\):
\[
    \operatorname{mult}_{\Schur_\lambda V}(M)
    =
    \dim\{v\in M_\lambda:\mathfrak n_+v=0\},
\]
where \(\mathfrak n_+\) is the nilpotent Lie algebra of the chosen positive
Borel.

The Lie algebra \(\mathfrak n_+\) is generated by the simple positive-root
operators
\[
    E_{1,2},E_{2,3},\ldots,E_{8,9}.
\]
Therefore a vector \(v\in M_\lambda\) is killed by all of \(\mathfrak n_+\) if
and only if it is killed by each simple raising operator \(E_{i,i+1}\). Hence
\[
    \{v\in M_\lambda:\mathfrak n_+v=0\}
    =
    \bigcap_{i=1}^{8}
    \ker\left(
        E_{i,i+1}:M_\lambda\to
        M_{\lambda+\varepsilon_i-\varepsilon_{i+1}}
    \right).
\]
This proves the formula.
\end{proof}

\begin{corollary}[Plethysm, highest weights, and invariants agree]
Let \(M_d=\Sym^d(\Sym^3V)\). If \(3d=9k\), then
\[
    \dim M_d^{SL_9}
    =
    \big\langle h_d[h_3],s_{(k^9)}\big\rangle
    =
    \dim
    \bigcap_{i=1}^{8}
    \ker\left(
        E_{i,i+1}:(M_d)_{(k^9)}
        \to
        (M_d)_{(k^9)+\varepsilon_i-\varepsilon_{i+1}}
    \right).
\]
If \(3d\) is not divisible by \(9\), then
\[
    \dim M_d^{SL_9}=0.
\]
\end{corollary}

\begin{proof}
The first equality is the invariant-dimension theorem. The second equality
follows from the highest-weight kernel proposition applied to
\(M=M_d\) and \(\lambda=(k^9)\), because the multiplicity of
\(\Schur_{(k^9)}V\) is the Schur coefficient
\[
    \big\langle h_d[h_3],s_{(k^9)}\big\rangle.
\]
The zero statement in the non-divisible case follows from the same theorem.
\end{proof}

\section*{4. Consequence for the certified degrees}

The theorem above is purely representation-theoretic. It does not use any
numerical vanishing of a coefficient in its proof.

For \(d=3,6,9\), the possible rectangular weights are respectively
\[
    (1^9),\qquad (2^9),\qquad (3^9).
\]
Therefore the theorem gives
\[
    \dim \Sym^3(\Sym^3\C^9)^{SL_9}
    =
    \big\langle h_3[h_3],s_{(1^9)}\big\rangle,
\]
\[
    \dim \Sym^6(\Sym^3\C^9)^{SL_9}
    =
    \big\langle h_6[h_3],s_{(2^9)}\big\rangle,
\]
and
\[
    \dim \Sym^9(\Sym^3\C^9)^{SL_9}
    =
    \big\langle h_9[h_3],s_{(3^9)}\big\rangle.
\]
Combining this Proved-Alg equivalence with exact independent character
computations of these three coefficients gives
\[
    \dim \Sym^d(\Sym^3\C^9)^{SL_9}=0
    \qquad
    \text{for }d\in\{3,6,9\}.
\]

\begin{remark}[Separation of proof and computation]
The equality
\[
    \dim \Sym^d(\Sym^3\C^9)^{SL_9}
    =
    \big\langle h_d[h_3],s_{(k^9)}\big\rangle
\]
is Proved-Alg. The numerical vanishing of the coefficients in the listed
degrees is a separate exact character computation. A direct highest-weight rank
computation is not required for the proof of the equality; when feasible, it is
only an independent way to compute the same multiplicity.
\end{remark}

\end{document}
